\renewcommand{\thesection}{\arabic{section}}\renewcommand{\theHsection}{chapter.\arabic{section}}\setcounter{section}{9}\section{多变量函数的重积分}
\noindent\begin{minipage}{\textwidth}\ProblemHeading{例题}{10.1}
计算 $\displaystyle\iint_D e^{x+y}\,dx\,dy$，其中 $D=[0,1]\times[0,1]$。
\end{minipage}\par\vfill
\noindent\begin{minipage}{\textwidth}\ProblemHeading{例题}{10.2}
计算 $\displaystyle\iint_D x\cos(xy)\,dx\,dy$，其中 $D=[0,\pi]\times[0,1]$。
\end{minipage}\par\vfill
\noindent\begin{minipage}{\textwidth}\ProblemHeading{例题}{10.3}
设 $a<b$，$f(x,y)$ 在 $[a,b]\times[a,b]$ 上连续。证明\[\int_a^b dx\int_a^x f(x,y)\,dy=\int_a^b dy\int_y^b f(x,y)\,dx.\]
\end{minipage}\par\vfill
\noindent\begin{minipage}{\textwidth}\ProblemHeading{例题}{10.4}
计算 $\displaystyle\iint_D x^2y^2\,dx\,dy$，其中 $D$ 由 $y=0$、$x=a$ 和 $y=(b/a)x$ 围成，$a,b>0$。
\end{minipage}\par\vfill
\newpage
\noindent\begin{minipage}{\textwidth}\ProblemHeading{例题}{10.5}
计算 $\displaystyle\iint_D(x^2+y^2)\,dx\,dy$，其中 $D$ 是由 $y=a$、$y=3a$、$y=x$、$y=x+a$ 围成的平行四边形，$a>0$。
\end{minipage}\par\vfill
\noindent\begin{minipage}{\textwidth}\ProblemHeading{例题}{10.6}
计算 $\displaystyle\int_1^2 y\,dy\int_y^2\frac{\sin x}{x^2-1}\,dx$（在端点 $y=1$ 处按反常积分理解）。
\end{minipage}\par\vfill
\noindent\begin{minipage}{\textwidth}\ProblemHeading{例题}{10.7}
计算 $\displaystyle\iint_{|x|+|y|\le1}x^2y^2\,dx\,dy$。
\end{minipage}\par\vfill
\noindent\begin{minipage}{\textwidth}\ProblemHeading{例题}{10.8}
求由 $x^2+y^2\le a^2$ 和 $x^2+z^2\le a^2$ 相交得到的立体的体积，其中 $a>0$。
\end{minipage}\par\vfill
\newpage
\noindent\begin{minipage}{\textwidth}\ProblemHeading{习题}{10.9}
改变下列累次积分的积分次序（设所涉及的函数连续，参数 $a>0$、$a<b$ 分别按各式需要取值）：
\SubQ{1}{$\int_{-1}^{1}dx\int_0^{\sqrt{1-x^2}} f(x,y)\,dy$。}
\SubQ{2}{$\int_0^2dx\int_{2x}^{6-x}f(x,y)\,dy$。}
\SubQ{3}{$\int_0^ady\int_{a-\sqrt{a^2-y^2}}^{a+\sqrt{a^2-y^2}}f(x,y)\,dx$。}
\end{minipage}\par\vfill
\noindent\begin{minipage}{\textwidth}\ProblemHeading{习题}{10.10}
改变下列累次积分的积分次序（设所涉及的函数连续，参数 $a>0$、$a<b$ 分别按各式需要取值）：
\SubQ{1}{$\int_a^bdy\int_y^bf(x,y)\,dx$。}
\SubQ{2}{$\int_0^1dx\int_0^xf(x,y)\,dy+\int_1^2dx\int_0^{2-x}f(x,y)\,dy$。}
\SubQ{3}{$\int_0^1dy\int_{1/2}^1f(x,y)\,dx+\int_1^2dy\int_{1/2}^{1/y}f(x,y)\,dx$。}
\end{minipage}\par\vfill
\noindent\begin{minipage}{\textwidth}\ProblemHeading{习题}{10.11}
计算下列二重积分，出现的 $a$ 取正数：
\SubQ{1}{$\iint_D\frac{y}{(1+x^2+y^2)^{3/2}}\,dx\,dy$，$D=[0,1]\times[0,1]$。}
\SubQ{2}{$\iint_D\sin(x+y)\,dx\,dy$，$D=[0,\pi]\times[0,\pi]$。}
\SubQ{3}{$\iint_D\cos(x+y)\,dx\,dy$，$D$ 由 $y=\pi$、$x=y$、$x=0$ 围成。}
\end{minipage}\par\vfill
\noindent\begin{minipage}{\textwidth}\ProblemHeading{习题}{10.12}
计算下列二重积分，出现的 $a$ 取正数：
\SubQ{1}{$\iint_D(x+y)\,dx\,dy$，$D$ 是圆 $x^2+y^2\le a^2$ 的第一象限部分。}
\SubQ{2}{$\iint_D(x+y-1)\,dx\,dy$，$D$ 由 $y=x$、$y=x+a$、$y=a$、$y=3a$ 围成。}
\SubQ{3}{$\iint_D\frac{\sin y}{y}\,dx\,dy$，$D$ 由 $y=x$ 和 $x=y^2$ 围成；被积函数在 $y=0$ 处按连续延拓取值。}
\end{minipage}\par\vfill
\newpage
\noindent\begin{minipage}{\textwidth}\ProblemHeading{习题}{10.13}
计算下列二重积分，出现的 $a$ 取正数：
\SubQ{1}{$\iint_D\frac{x^2}{y^2}\,dx\,dy$，$D$ 由 $x=2$、$y=x$ 和 $xy=1$ 围成。}
\SubQ{2}{$\iint_D|\cos(x+y)|\,dx\,dy$，$D$ 由 $y=x$、$y=0$ 和 $x=\pi/2$ 围成。}
\end{minipage}\par\vfill
\noindent\begin{minipage}{\textwidth}\ProblemHeading{习题}{10.14}
利用函数的奇偶性计算下列积分：
\SubQ{1}{$\iint_D(x^2+y^2)\,dx\,dy$，$D=[-1,1]\times[-1,1]$。}
\SubQ{2}{$\iint_D\sin x\sin y\,dx\,dy$，$D$ 是由 $x^2-y^2=1$、$x^2+y^2=9$ 围成且含原点的部分。}
\end{minipage}\par\vfill
\noindent\begin{minipage}{\textwidth}\ProblemHeading{习题}{10.15}
设 $\varphi$ 和 $\psi$ 分别在 $[a,b]$ 和 $[c,d]$ 上可积。证明 $f(x,y)=\varphi(x)\psi(y)$ 在 $D=[a,b]\times[c,d]$ 上可积，且\[\iint_D f(x,y)\,dx\,dy=\left(\int_a^b\varphi(x)\,dx\right)\left(\int_c^d\psi(y)\,dy\right).\]
\end{minipage}\par\vfill
\noindent\begin{minipage}{\textwidth}\ProblemHeading{习题}{10.16}
设 $a>0$，$f$ 在 $[0,a]$ 上连续。证明：
\SubQ{1}{$\int_0^a dx\int_0^x f(x)f(y)\,dy=\frac12\left(\int_0^a f(x)\,dx\right)^2$。}
\SubQ{2}{$\int_0^a dx\int_0^x f(y)\,dy=\int_0^a(a-x)f(x)\,dx$。}
\end{minipage}\par\vfill
\newpage
\noindent\begin{minipage}{\textwidth}\ProblemHeading{习题}{10.17}
设 $f(x,y)$ 有连续的二阶偏导数，$D=[a,b]\times[c,d]$。求 $\displaystyle\iint_D\frac{\partial^2f}{\partial x\partial y}\,dx\,dy$。
\end{minipage}\par\vfill
\noindent\begin{minipage}{\textwidth}\ProblemHeading{习题}{10.18}
设 $f(x,y)$ 在平面上连续。求极限 $\displaystyle\lim_{r\to0^+}\frac1{\pi r^2}\iint_{x^2+y^2\le r^2}f(x,y)\,dx\,dy$。
\end{minipage}\par\vfill
\noindent\begin{minipage}{\textwidth}\ProblemHeading{例题}{10.19（改写）}
计算 $\displaystyle\iint_D\sqrt{x^2+y^2}\,dx\,dy$，其中 $D=\{(x,y):x^2+y^2\le1,\ x\ge0,\ y\ge0\}$。
\end{minipage}\par\vfill
\noindent\begin{minipage}{\textwidth}\ProblemHeading{例题}{10.20}
设 $a,b,c>0$。求椭球 $\displaystyle\frac{x^2}{a^2}+\frac{y^2}{b^2}+\frac{z^2}{c^2}\le1$ 的体积。
\end{minipage}\par\vfill
\newpage
\noindent\begin{minipage}{\textwidth}\ProblemHeading{例题}{10.21}
设 $a>0$。求球 $x^2+y^2+z^2\le a^2$ 在圆柱 $x^2+y^2\le ay$ 内的部分的体积。
\end{minipage}\par\vfill
\noindent\begin{minipage}{\textwidth}\ProblemHeading{例题}{10.22}
设 $R>0$。求两个球 $x^2+y^2+z^2\le R^2$ 和 $x^2+y^2+(z-R)^2\le R^2$ 的公共部分的体积。
\end{minipage}\par\vfill
\noindent\begin{minipage}{\textwidth}\ProblemHeading{例题}{10.23}
设 $a>0$。求双纽线 $(x^2+y^2)^2=2a^2(x^2-y^2)$ 所围成的面积。
\end{minipage}\par\vfill
\noindent\begin{minipage}{\textwidth}\ProblemHeading{例题}{10.24}
求 $\displaystyle\iint_D(\sqrt{x}+\sqrt{y})\,dx\,dy$，其中 $D$ 是由 $\sqrt{x}+\sqrt{y}=1$ 与两坐标轴围成的区域。
\end{minipage}\par\vfill
\newpage
\noindent\begin{minipage}{\textwidth}\ProblemHeading{例题}{10.25}
求 $\displaystyle\iint_Dx^2y^2\,dx\,dy$，其中 $D$ 由四条抛物线 $y^2=px$、$y^2=qx$、$x^2=ay$ 和 $x^2=by$ 围成，$0<p<q$、$0<a<b$。
\end{minipage}\par\vfill
\noindent\begin{minipage}{\textwidth}\ProblemHeading{例题}{10.26（改写）}
设 $R>0$。计算下列积分：
\SubQ{1}{$\displaystyle\iint_{x^2+y^2\le R^2}e^{-x^2-y^2}\,dx\,dy$。}
\SubQ{2}{$\displaystyle\int_{-\infty}^{+\infty}e^{-x^2}\,dx$。}
\end{minipage}\par\vfill
\noindent\begin{minipage}{\textwidth}\ProblemHeading{习题}{10.27}
计算下列积分，出现的 $R,a,b$ 均为正数：
\SubQ{1}{$\int_0^Rdx\int_0^{\sqrt{R^2-x^2}}\ln(1+x^2+y^2)\,dy$。}
\SubQ{2}{$\int_0^adx\int_0^bxy(x^2-y^2)\,dy$。}
\SubQ{3}{$\int_0^\pi\int_0^\pi\cos(x+y)\,dx\,dy$。}
\end{minipage}\par\vfill
\noindent\begin{minipage}{\textwidth}\ProblemHeading{习题}{10.28}
计算下列积分，出现的 $R,a,b$ 均为正数：
\SubQ{1}{$\int_0^{1/\sqrt2}dx\int_x^{\sqrt{1-x^2}}xy(x+y)\,dy$。}
\SubQ{2}{$\int_0^{R/\sqrt{1+R^2}}dx\int_0^{Rx}\left(1+\frac{y^2}{x^2}\right)dy+\int_{R/\sqrt{1+R^2}}^Rdx\int_0^{\sqrt{R^2-x^2}}\left(1+\frac{y^2}{x^2}\right)dy$。}
\end{minipage}\par\vfill
\newpage
\noindent\begin{minipage}{\textwidth}\ProblemHeading{习题}{10.29}
计算下列二重积分；未另注明时，出现的长度参数均取正数：
\SubQ{1}{$\iint_D\sqrt{x^2+y^2}\,dx\,dy$，$D:\ x^2+y^2\le x+y$。}
\SubQ{2}{$\iint_D\sqrt{\frac{x^2}{a^2}+\frac{y^2}{b^2}}\,dx\,dy$，$D$ 为 $x^2/a^2+y^2/b^2=4$、$y=0$、$y=x$ 围成的第一象限部分。}
\SubQ{3}{$\iint_D(x^2+y^2)\,dx\,dy$，$D$ 为 $xy=1$、$xy=2$、$y=x$、$y=2x$ 围成的第一象限部分。}
\end{minipage}\par\vfill
\noindent\begin{minipage}{\textwidth}\ProblemHeading{习题}{10.30}
计算下列二重积分；未另注明时，出现的长度参数均取正数：
\SubQ{1}{$\iint_Ddx\,dy$，$D$ 由 $y^2=ax$、$y^2=bx$、$x^2=my$、$x^2=ny$ 围成，其中 $a>b>0$、$m>n>0$。}
\SubQ{2}{$\iint_Dxy\,dx\,dy$，$D$ 由 $xy=a$、$xy=b$、$y^2=cx$、$y^2=dx$ 围成的第一象限部分，其中 $0<a<b$、$0<c<d$。}
\SubQ{3}{$\iint_D4xy\,dx\,dy$，$D:\ x^4+y^4\le1,\ x\ge0,\ y\ge0$。}
\end{minipage}\par\vfill
\noindent\begin{minipage}{\textwidth}\ProblemHeading{习题}{10.31}
计算下列二重积分；未另注明时，出现的长度参数均取正数：
\SubQ{1}{$\iint_D\frac{x^2-y^2}{\sqrt{x+y+3}}\,dx\,dy$，$D:\ |x|+|y|\le1$。}
\SubQ{2}{$\iint_D\sin\frac{y}{x+y}\,dx\,dy$，$D$ 由 $x+y=1$、$x=0$、$y=0$ 围成；原点处的值可任意定义。}
\SubQ{3}{$\iint_D|xy|\,dx\,dy$，$D:\ x^2+y^2\le a^2$。}
\end{minipage}\par\vfill
\noindent\begin{minipage}{\textwidth}\ProblemHeading{习题}{10.32}
求下列曲线围成的平面区域的面积：
\SubQ{1}{$x^2+2y^2=3$ 和 $xy=1$ 围成的不含原点部分。}
\SubQ{2}{$(x-y)^2+x^2=a^2$，其中 $a>0$。}
\SubQ{3}{$x+y=a$、$x+y=b$、$y=kx$ 和 $y=mx$，其中 $0<a<b$、$0<k<m$。}
\end{minipage}\par\vfill
\newpage
\noindent\begin{minipage}{\textwidth}\ProblemHeading{习题}{10.33}
证明 $\displaystyle\iint_{x^2+y^2\le1}e^{x^2+y^2}\,dx\,dy\le\left(\int_{-\sqrt{\pi}/2}^{\sqrt{\pi}/2}e^{x^2}\,dx\right)^2$。
\end{minipage}\par\vfill
\noindent\begin{minipage}{\textwidth}\ProblemHeading{习题}{10.34}
设 $f$ 在 $[0,1]$ 上连续。证明 $\displaystyle\left(\int_0^1e^{f(x)}\,dx\right)\left(\int_0^1e^{-f(y)}\,dy\right)\ge1$。
\end{minipage}\par\vfill
\noindent\begin{minipage}{\textwidth}\ProblemHeading{习题}{10.35}
设 $f$ 为连续的奇函数。证明 $\displaystyle\iint_{|x|+|y|\le1}e^{f(x+y)}\,dx\,dy\ge2$。
\end{minipage}\par\vfill
\noindent\begin{minipage}{\textwidth}\ProblemHeading{习题}{10.36}
设 $A>0$，$f$ 为连续函数，$D=\{(x,y):|x|\le A/2,\ |y|\le A/2\}$。证明\[\iint_Df(x-y)\,dx\,dy=\int_{-A}^{A}f(t)(A-|t|)\,dt.\]
\end{minipage}\par\vfill
\newpage
\noindent\begin{minipage}{\textwidth}\ProblemHeading{例题}{10.37}
求 $\displaystyle\iiint_Vx^2ye^{xyz}\,dx\,dy\,dz$，其中 $V=[0,1]^3$。
\end{minipage}\par\vfill
\noindent\begin{minipage}{\textwidth}\ProblemHeading{例题}{10.38}
计算 $\displaystyle\iiint_V\frac{dx\,dy\,dz}{(1+x+y+z)^3}$，其中 $V$ 是由三个坐标平面与平面 $x+y+z=1$ 围成的四面体。
\end{minipage}\par\vfill
\noindent\begin{minipage}{\textwidth}\ProblemHeading{例题}{10.39}
设 $R,h>0$。计算 $\displaystyle\iiint_Vz\,dx\,dy\,dz$，其中 $V$ 是由锥面 $R^2z^2=h^2(x^2+y^2)$ 与平面 $z=h$ 围成的有界锥体。
\end{minipage}\par\vfill
\noindent\begin{minipage}{\textwidth}\ProblemHeading{例题}{10.40}
设 $a,b,c>0$。计算 $\displaystyle\iiint_V(x^2+y^2+z^2)\,dx\,dy\,dz$，其中 $V:\ x^2/a^2+y^2/b^2+z^2/c^2\le1$。
\end{minipage}\par\vfill
\newpage
\noindent\begin{minipage}{\textwidth}\ProblemHeading{例题}{10.41}
设 $a>0$。求 $\displaystyle\iiint_Vz\sqrt{x^2+y^2}\,dx\,dy\,dz$，其中 $V$ 由柱面 $x^2+y^2=2x$ 与平面 $z=0$、$z=a$ 围成。
\end{minipage}\par\vfill
\noindent\begin{minipage}{\textwidth}\ProblemHeading{例题}{10.42}
设 $h>0$。求 $\displaystyle\iiint_V\frac{dx\,dy\,dz}{1+x^2+y^2}$，其中 $V:\ 0\le z\le h,\ x^2+y^2\le4z$。
\end{minipage}\par\vfill
\noindent\begin{minipage}{\textwidth}\ProblemHeading{例题}{10.43}
设 $R>0$。计算 $\displaystyle\iiint_V(x^2+y^2+z^2)\,dx\,dy\,dz$，其中 $V$ 是球 $x^2+y^2+z^2\le R^2$ 中位于锥面 $z=\sqrt{x^2+y^2}$ 上方的部分。
\end{minipage}\par\vfill
\noindent\begin{minipage}{\textwidth}\ProblemHeading{例题}{10.44}
设 $a>0$。求曲面 $(x^2+y^2+z^2)^2=a^3z$ 所围成的立体的体积。
\end{minipage}\par\vfill
\newpage
\noindent\begin{minipage}{\textwidth}\ProblemHeading{例题}{10.45}
设 $a,b,c>0$。求曲面 $\displaystyle\left(\frac xa\right)^{2/3}+\left(\frac yb\right)^{2/3}+\left(\frac zc\right)^{2/3}=1$ 所围成立体的体积，其中 $t^{2/3}=(\sqrt[3]{t})^2$。
\end{minipage}\par\vfill
\noindent\begin{minipage}{\textwidth}\ProblemHeading{例题}{10.46}
设 $a>0$，球体 $x^2+y^2+z^2\le2az$ 内每一点的密度与该点到原点的距离成正比。求此球体的质心。
\end{minipage}\par\vfill
\noindent\begin{minipage}{\textwidth}\ProblemHeading{例题}{10.47}
求半径为 $R$、高为 $l$、密度为常数 $\rho$ 的均匀圆柱体对其轴线的转动惯量。
\end{minipage}\par\vfill
\noindent\begin{minipage}{\textwidth}\ProblemHeading{例题}{10.48}
求半径为 $R$、质量为 $M$ 的均匀球体绕其直径的转动惯量。
\end{minipage}\par\vfill
\newpage
\noindent\begin{minipage}{\textwidth}\ProblemHeading{例题}{10.49}
求半径为 $R$、密度为常数 $\rho$ 的均匀球体对球外质量为 $m$ 的质点 $A$ 的万有引力。设 $A$ 到球心的距离为 $l>R$，万有引力常数为 $G$。
\end{minipage}\par\vfill
\noindent\begin{minipage}{\textwidth}\ProblemHeading{习题}{10.50}
计算下列三重积分，出现的参数 $a>0$：
\SubQ{1}{$\iiint_Vxy\,dx\,dy\,dz$，$V:\ 1\le x\le2,\ -2\le y\le1,\ 0\le z\le1/2$。}
\SubQ{2}{$\iiint_Vxy^2z^3\,dx\,dy\,dz$，$V$ 由 $z=xy$、$y=x$、$x=1$、$z=0$ 围成的有界区域。}
\SubQ{3}{$\iiint_Vy\cos(x+z)\,dx\,dy\,dz$，$V$ 由 $y=\sqrt{x}$、$y=0$、$z=0$ 和 $x+z=\pi/2$ 围成。}
\end{minipage}\par\vfill
\noindent\begin{minipage}{\textwidth}\ProblemHeading{习题}{10.51}
计算下列三重积分，出现的参数 $a>0$：
\SubQ{1}{$\iiint_V(a-y)\,dx\,dy\,dz$，$V$ 由 $y=0$、$z=0$、$2x+y=a$、$x+y=a$ 和 $y+z=a$ 围成。}
\end{minipage}\par\vfill
\noindent\begin{minipage}{\textwidth}\ProblemHeading{习题}{10.52}
计算下列累次积分，$a,R>0$：
\SubQ{1}{$\int_0^2dx\int_0^{\sqrt{2x-x^2}}dy\int_0^az\sqrt{x^2+y^2}\,dz$。}
\SubQ{2}{$\int_{-R}^Rdx\int_{-\sqrt{R^2-x^2}}^{\sqrt{R^2-x^2}}dy\int_0^{\sqrt{R^2-x^2-y^2}}(x^2+y^2)\,dz$。}
\SubQ{3}{$\int_0^1dx\int_0^{\sqrt{1-x^2}}dy\int_0^{\sqrt{1-x^2-y^2}}\sqrt{x^2+y^2+z^2}\,dz$。}
\end{minipage}\par\vfill
\newpage
\noindent\begin{minipage}{\textwidth}\ProblemHeading{习题}{10.53}
计算下列累次积分，$a,R>0$：
\SubQ{1}{$\int_0^1dx\int_0^{\sqrt{1-x^2}}dy\int_{\sqrt{x^2+y^2}}^{\sqrt{2-x^2-y^2}}z^2\,dz$。}
\end{minipage}\par\vfill
\noindent\begin{minipage}{\textwidth}\ProblemHeading{习题}{10.54}
计算下列三重积分：
\SubQ{1}{$\iiint_V(x^2+y^2)\,dx\,dy\,dz$，$V$ 由 $x^2+y^2=2z$、$z=2$ 围成。}
\SubQ{2}{$\iiint_V\sqrt{x^2+y^2}\,dx\,dy\,dz$，$V$ 由 $x^2+y^2=z^2$、$z=1$ 围成的有界区域。}
\SubQ{3}{$\iiint_Vz\,dx\,dy\,dz$，$V$ 由 $z=\sqrt{4-x^2-y^2}$、$x^2+y^2=3z$ 围成。}
\end{minipage}\par\vfill
\noindent\begin{minipage}{\textwidth}\ProblemHeading{习题}{10.55}
计算下列三重积分：
\SubQ{1}{$\iiint_Vxyz\,dx\,dy\,dz$，$V$ 是单位球的第一卦限部分。}
\SubQ{2}{$\iiint_Vx^2\,dx\,dy\,dz$，$V$ 由 $z=y^2$、$z=4y^2$（$y>0$）以及 $z=x$、$z=2x$、$z=1$ 围成。}
\SubQ{3}{$\iiint_V|x^2+y^2+z^2-1|\,dx\,dy\,dz$，$V:\ x^2+y^2+z^2\le4$。}
\end{minipage}\par\vfill
\noindent\begin{minipage}{\textwidth}\ProblemHeading{习题}{10.56}
计算下列三重积分：
\SubQ{1}{$\iiint_Ve^{|z|}\,dx\,dy\,dz$，$V:\ x^2+y^2+z^2\le1$。}
\SubQ{2}{$\iiint_V(|x|+z)e^{-(x^2+y^2+z^2)}\,dx\,dy\,dz$，$V:\ 1\le x^2+y^2+z^2\le4$。}
\end{minipage}\par\vfill
\newpage
\noindent\begin{minipage}{\textwidth}\ProblemHeading{习题}{10.57}
利用对称性计算下列三重积分，参数 $a,b,c>0$、$0\le r<R$：
\SubQ{1}{$\iiint_V(x+y)\,dx\,dy\,dz$，$V$ 由 $z=1-x^2-y^2$、$z=0$ 围成。}
\SubQ{2}{$\iiint_Vx\,dx\,dy\,dz$，$V$ 由 $x^2+y^2=z^2$ 与 $x^2+y^2=1$ 围成的有界区域。}
\SubQ{3}{$\iiint_V\sqrt{x^2+y^2+z^2}\,dx\,dy\,dz$，$V:\ x^2+y^2+z^2\le x$。}
\end{minipage}\par\vfill
\noindent\begin{minipage}{\textwidth}\ProblemHeading{习题}{10.58}
利用对称性计算下列三重积分，参数 $a,b,c>0$、$0\le r<R$：
\SubQ{1}{$\iiint_V(x^2+y^2)\,dx\,dy\,dz$，$V:\ r^2\le x^2+y^2+z^2\le R^2,\ z\ge0$。}
\SubQ{2}{$\iiint_V\sqrt{1-x^2/a^2-y^2/b^2-z^2/c^2}\,dx\,dy\,dz$，$V:\ x^2/a^2+y^2/b^2+z^2/c^2\le1$。}
\SubQ{3}{$\iiint_V\frac{z\ln(x^2+y^2+z^2+1)}{x^2+y^2+z^2+1}\,dx\,dy\,dz$，$V:\ x^2+y^2+z^2\le1$。}
\end{minipage}\par\vfill
\noindent\begin{minipage}{\textwidth}\ProblemHeading{习题}{10.59}
计算下列曲面围成立体的体积，参数 $a,b,c>0$：
\SubQ{1}{$y=0$、$z=0$、$3x+y=6$、$3x+2y=12$、$x+y+z=6$。}
\SubQ{2}{$z=x^2+y^2$、$z=2x^2+2y^2$、$y=x$、$y=x^2$。}
\SubQ{3}{$z^2+x^2=1$、$x+y+z=3$、$y=0$。}
\end{minipage}\par\vfill
\noindent\begin{minipage}{\textwidth}\ProblemHeading{习题}{10.60}
计算下列曲面围成立体的体积，参数 $a,b,c>0$：
\SubQ{1}{$x^2+y^2=2x$、$z=x^2+y^2$、$z=0$。}
\SubQ{2}{$x^2/9+y^2/4=1$、$z=xy$ 围成的第一卦限部分。}
\SubQ{3}{$x^2+y^2+z^2=2az$、$x^2+y^2=z^2$ 围成的含 $z$ 轴部分。}
\end{minipage}\par\vfill
\newpage
\noindent\begin{minipage}{\textwidth}\ProblemHeading{习题}{10.61}
计算下列曲面围成立体的体积，参数 $a,b,c>0$：
\SubQ{1}{$x^2/a^2+y^2/b^2+z^2/c^2=1$、$x^2/a^2+y^2/b^2=z^2/c^2$ 围成的含 $z$ 轴部分。}
\SubQ{2}{$(x^2+y^2+z^2)^2=a^3x$。}
\end{minipage}\par\vfill
\noindent\begin{minipage}{\textwidth}\ProblemHeading{习题}{10.62}
求函数 $f(x,y,z)=x^2+y^2+z^2$ 在区域 $x^2+y^2+z^2\le x+y+z$ 内的平均值。
\end{minipage}\par\vfill
\noindent\begin{minipage}{\textwidth}\ProblemHeading{习题}{10.63}
设 $f$ 为可微函数，$\displaystyle F(t)=\iiint_{x^2+y^2+z^2\le t^2}f(x^2+y^2+z^2)\,dx\,dy\,dz$。求 $F^{\prime}(t)$。
\end{minipage}\par\vfill
\noindent\begin{minipage}{\textwidth}\ProblemHeading{习题}{10.64}
设 $f$ 在 $[-1,1]$ 上连续。证明 $\displaystyle\iiint_{x^2+y^2+z^2\le1}f(z)\,dx\,dy\,dz=\pi\int_{-1}^{1}f(z)(1-z^2)\,dz$。
\end{minipage}\par\vfill
\newpage
\noindent\begin{minipage}{\textwidth}\ProblemHeading{习题}{10.65}
设 $a<b$，$f(x,y,z)$ 连续。证明\[\int_a^bdx\int_a^xdy\int_a^yf(x,y,z)\,dz=\int_a^bdz\int_z^bdy\int_y^bf(x,y,z)\,dx.\]
\end{minipage}\par\vfill
\noindent\begin{minipage}{\textwidth}\ProblemHeading{习题}{10.66}
设 $a,b>0$。求椭圆薄片 $x^2/a^2+y^2/b^2\le1$ 的质量，密度为 $\rho=x^2/a^2+y^2/b^2$。
\end{minipage}\par\vfill
\noindent\begin{minipage}{\textwidth}\ProblemHeading{习题}{10.67}
平面圆环由半径 $R>r>0$ 的两个同心圆围成。各点面密度与到圆心的距离成反比，并且内圆周上的面密度为1。求圆环的质量。
\end{minipage}\par\vfill
\noindent\begin{minipage}{\textwidth}\ProblemHeading{习题}{10.68}
球壳由半径 $R>r>0$ 的两个同心球围成。材料各点的密度与到球心的距离成反比，并且该比例规律在距离1处给出密度 $k>0$。求物体的总质量。
\end{minipage}\par\vfill
\newpage
\noindent\begin{minipage}{\textwidth}\ProblemHeading{习题}{10.69}
半径为 $a>0$ 的圆盘，各点面密度等于到圆心的距离。截去一个与大圆内切、半径为 $a/2$ 的小圆盘。取大圆心为原点，小圆心在 $x$ 轴正半轴上，求剩余部分的质心坐标。
\end{minipage}\par\vfill
\noindent\begin{minipage}{\textwidth}\ProblemHeading{习题}{10.70}
均质薄板由半径 $a>0$ 的半圆和一个长方形沿半圆直径拼接而成；长方形长为 $2a$，宽为 $b$。为使薄板质心恰在半圆圆心，$b$ 应取多少？
\end{minipage}\par\vfill
\noindent\begin{minipage}{\textwidth}\ProblemHeading{习题}{10.71}
设 $a,b,c>0$。求由锥面 $x^2/a^2+y^2/b^2=z^2/c^2$ 与平面 $z=c$ 围成的均匀物体的质心坐标。
\end{minipage}\par\vfill
\noindent\begin{minipage}{\textwidth}\ProblemHeading{习题}{10.72}
设 $a>0$。球体 $x^2+y^2+z^2\le2az$ 内各点密度与到原点的距离成反比。求其质心坐标。
\end{minipage}\par\vfill
\newpage
\noindent\begin{minipage}{\textwidth}\ProblemHeading{习题}{10.73}
在半径为 $a>0$ 的圆柱顶面连接一个半径为 $a$ 的半球，二者密度相同且均匀。为使整体质心恰在半球的球心处，圆柱高应为多少？
\end{minipage}\par\vfill
\noindent\begin{minipage}{\textwidth}\ProblemHeading{习题}{10.74}
求下列均匀物体的转动惯量：
\SubQ{1}{质量为 $m$、半径为 $R$ 的薄圆盘，对通过圆心并垂直于圆盘的轴。}
\SubQ{2}{质量为 $m$、半径为 $R$ 的薄圆盘，对其一条直径。}
\SubQ{3}{质量为 $m$、半径为 $R$ 的球体，对通过球心的轴。}
\end{minipage}\par\vfill
\noindent\begin{minipage}{\textwidth}\ProblemHeading{习题}{10.75}
求密度为常数 $\rho$ 的均匀球锥体对位于其顶点的单位质量质点的万有引力。球的半径为 $R$，球锥体的轴截面为圆心角 $2\alpha$ 的扇形，$0<\alpha<\pi$。以 $G$ 表示万有引力常数。
\end{minipage}\par\vfill
\noindent\begin{minipage}{\textwidth}\ProblemHeading{例题}{10.76（改写）}
设 $a>0$，$n$ 为正整数，$S_n(a)=\{(x_1,\ldots,x_n)\in\mathbb R^n:x_i\ge0,\ \sum_{i=1}^nx_i\le a\}$。求 $n$ 维单形 $S_n(a)$ 的体积，并推导关于维数的递推关系。
\end{minipage}\par\vfill
\newpage
\noindent\begin{minipage}{\textwidth}\ProblemHeading{例题}{10.77（改写）}
设 $a>0$，$B_n(a)=\{x\in\mathbb R^n:\sum_{i=1}^nx_i^2\le a^2\}$。求 $n$ 维球体 $B_n(a)$ 的体积，并推导单位球体体积关于维数的递推关系。
\end{minipage}\par\vfill
\noindent\begin{minipage}{\textwidth}\ProblemHeading{习题}{10.78}
设 $n$ 为正整数。计算下列 $n$ 重积分：
\SubQ{1}{$\int_{[0,1]^n}(x_1^2+\cdots+x_n^2)\,dx_1\cdots dx_n$。}
\SubQ{2}{$\int_{[0,1]^n}(x_1+\cdots+x_n)^2\,dx_1\cdots dx_n$。}
\SubQ{3}{$\int_0^1dx_1\int_0^{x_1}dx_2\cdots\int_0^{x_{n-1}}x_1x_2\cdots x_n\,dx_n$。}
\end{minipage}\par\vfill
\noindent\begin{minipage}{\textwidth}\ProblemHeading{习题}{10.79}
设 $a_1,\ldots,a_n>0$。计算集合 $\displaystyle V_n=\left\{(x_1,\ldots,x_n):\sum_{i=1}^n\frac{x_i}{a_i}\le1,\ x_i\ge0\right\}$ 的 $n$ 维体积。
\end{minipage}\par\vfill
\noindent\begin{minipage}{\textwidth}\ProblemHeading{习题}{10.80}
设 $a>0$，$n$ 为正整数，$f$ 在 $[0,a]$ 上连续。证明\[\int_0^adx_1\int_0^{x_1}dx_2\cdots\int_0^{x_{n-1}}f(x_n)\,dx_n=\frac1{(n-1)!}\int_0^af(t)(a-t)^{n-1}\,dt.\]
\end{minipage}\par\vfill
\newpage
\noindent\begin{minipage}{\textwidth}\ProblemHeading{习题}{10.81}
设 $a>0$，$n$ 为正整数，$f$ 在 $[0,a]$ 上连续。证明\[\int_0^adx_1\int_0^{x_1}dx_2\cdots\int_0^{x_{n-1}}f(x_1)\cdots f(x_n)\,dx_n=\frac1{n!}\left(\int_0^af(t)\,dt\right)^n.\]
\end{minipage}\par\vfill
\noindent\begin{minipage}{\textwidth}\ProblemHeading{习题}{10.82}
计算 $\displaystyle I=\iint_D\operatorname{sgn}(x^2-y^2+2)\,dx\,dy$，其中 $D=\{(x,y):x^2+y^2\le4\}$，$\operatorname{sgn}$ 为符号函数。
\end{minipage}\par\vfill
\noindent\begin{minipage}{\textwidth}\ProblemHeading{习题}{10.83}
计算 $\displaystyle I=\iiint_{[0,1]^3}\frac{du\,dv\,dw}{(1+u^2+v^2+w^2)^2}$。
\end{minipage}\par\vfill
\noindent\begin{minipage}{\textwidth}\ProblemHeading{习题}{10.84}
设 $a,b>0$。求下列积分：
\SubQ{1}{$I_1=\int_0^1\sin\!\left(\ln\frac1x\right)\frac{x^b-x^a}{\ln x}\,dx$。}
\SubQ{2}{$I_2=\int_0^1\cos\!\left(\ln\frac1x\right)\frac{x^b-x^a}{\ln x}\,dx$。}
\end{minipage}\par\vfill
\newpage
\noindent\begin{minipage}{\textwidth}\ProblemHeading{习题}{10.85}
设 $D=\{(x,y):x^2+y^2\le1\}$。求 $\displaystyle I=\iint_D\left|\frac{x+y}{\sqrt2}-x^2-y^2\right|\,dx\,dy$。
\end{minipage}\par\vfill
\noindent\begin{minipage}{\textwidth}\ProblemHeading{习题}{10.86}
设 $a>0$。求圆盘 $(x-a)^2+(y-a)^2\le a^2$ 与曲线 $(x^2+y^2)^2=8a^2xy$ 所围区域的公共部分的面积。
\end{minipage}\par\vfill
\noindent\begin{minipage}{\textwidth}\ProblemHeading{习题}{10.87}
计算曲面 $(x^2+y^2)^2+z^4=y$ 所围立体的体积。
\end{minipage}\par\vfill
\noindent\begin{minipage}{\textwidth}\ProblemHeading{习题}{10.88}
证明 $\displaystyle\iint_{[0,1]^2}(xy)^{xy}\,dx\,dy=\int_0^1t^t\,dt$；在零点处用连续延拓定义 $t^t$。
\end{minipage}\par\vfill
\newpage
\noindent\begin{minipage}{\textwidth}\ProblemHeading{习题}{10.89}
设常数 $a,b$ 不全为零，$c\in\mathbb R$，$f$ 在相应取值区间连续。证明\[\iint_{x^2+y^2\le1}f(ax+by+c)\,dx\,dy=2\int_{-1}^{1}\sqrt{1-t^2}\,f\!\left(t\sqrt{a^2+b^2}+c\right)dt.\]
\end{minipage}\par\vfill
\noindent\begin{minipage}{\textwidth}\ProblemHeading{习题}{10.90}
设 $f$ 为连续可导的单变量函数，$t>0$，令 $\displaystyle F(t)=\iint_{[0,t]^2}f(xy)\,dx\,dy$。证明：
\SubQ{1}{$F^{\prime}(t)=\frac2t\left[F(t)+\iint_{[0,t]^2}xyf^{\prime}(xy)\,dx\,dy\right]$。}
\SubQ{2}{$F^{\prime}(t)=\frac2t\int_0^{t^2}f(s)\,ds$。}
\end{minipage}\par\vfill
\noindent\begin{minipage}{\textwidth}\ProblemHeading{习题}{10.91}
（Poincaré 不等式）设 $a<b$，$\varphi,\psi$ 为 $[a,b]$ 上的连续函数且 $\varphi\le\psi$，$D=\{(x,y):a\le x\le b,\ \varphi(x)\le y\le\psi(x)\}$。设 $f$ 在 $D$ 上连续可微，并且 $f(x,\varphi(x))=0$。证明存在常数 $M>0$，使\[\iint_Df^2(x,y)\,dx\,dy\le M\iint_D\left(\frac{\partial f}{\partial y}(x,y)\right)^2dx\,dy.\]
\end{minipage}\par\vfill
\noindent\begin{minipage}{\textwidth}\ProblemHeading{习题}{10.92}
设 $a>0$、$n$ 为正整数，$\Omega_n(a)=\{x\in\mathbb R^n:x_i\ge0,\ \sum_{i=1}^nx_i\le a\}$。求积分 $\displaystyle I_n(a)=\int_{\Omega_n(a)}x_1x_2\cdots x_n\,dx_1\cdots dx_n$。
\end{minipage}\par\vfill
\newpage
\noindent\begin{minipage}{\textwidth}\ProblemHeading{习题}{10.93}
设 $a<b$，$n$ 为正整数，$f(x_1,\ldots,x_n)$ 连续。证明\[\int_a^bdx_1\int_a^{x_1}dx_2\cdots\int_a^{x_{n-1}}f(x_1,\ldots,x_n)\,dx_n=\int_a^bdx_n\int_{x_n}^bdx_{n-1}\cdots\int_{x_2}^bf(x_1,\ldots,x_n)\,dx_1.\]
\end{minipage}\par\vfill
